\section{A square-root lower bound with bounded witnesses}\label{sec:omega}
The upper bounds still allow large excursions. The following result gives a
literal finite search interval for the next excursion beyond any starting
point. It uses one zero of zeta and needs no assumption on the other zeros.

\begin{theorem}\label{thm:sqrt-lower}
For $A\ge1$, put
\[
 P=65536A,\qquad V=4096 AP,\qquad J=32VP,\qquad B=16(A+1)4^J.
\]
Some integer $n$ with $A\le n\le B$ satisfies
\[
                     32768^2 M(n)^2>n.
\]
The least such $n$ in this interval is a primitive recursive witness function
of $A$. Thus $|M(n)|/\sqrt n$ exceeds $1/32768$ arbitrarily far out.
\end{theorem}
\begin{proof}
We use the constructive zero certificate in the ProveIt packet on zeta
zeros\,\cite{zeropacket}. Its exact rational Euler--Maclaurin enclosures
certify a contraction of factor less than $1/64$ on the disk of radius
$10^{-6}$ centered at $1/2+14.1347251417347\,i$. Iteration gives a
primitive recursive name for the unique zero $\rho$ in this disk. The
functional equation and conjugation send this zero to $1-\bar\rho$ in
the same disk, so $\rho=1/2+i\gamma$ with $14<\gamma<15$. We import
this human-proved certificate and its functional-equation foundation.
The exact rational computation was rerun for this report and gave the same JSON data as the pinned repository file. The continuation $Z$ used below is the Euler summation formula at cutoff
one, the same one used to enclose the derivatives in that packet.

Write $a=1/32768$, $e=1/P$, $s=\rho+e$, and $z_0=s+2$.
Suppose that $|M(n)|\le a\sqrt n$ holds throughout the indicated finite
interval. Put $m_B(t)=M(\min(\lfloor t\rfloor,B))$ for $t\ge1$, and
$D_B(z)=\sum_{d=1}^{B}\mu(d)d^{-z}$. Then $|m_B(t)|\le A\sqrt t$
everywhere, using $|M(n)|\le n$ for $n<A$. For $t\ge A$, the stronger
$|m_B(t)|\le a\sqrt t$ holds, including beyond $B$ where the table is
constant. For $\Re z>0$, integration of each step gives
\[
        D_B(z)=z\int_1^\infty m_B(t)t^{-z-1}\,dt .
\]
Indeed the contribution of the jump $\mu(d)$ is
$\mu(d)z\int_d^\infty t^{-z-1}dt=\mu(d)d^{-z}$.
At $z=s$, integrate $|m_B(t)|\le t$ on $[1,A]$ and
$|m_B(t)|\le a\sqrt t$ on $[A,\infty)$. This gives
$|D_B(s)|\le16(2A+a/e)$ since $|s|\le16$.
The continuation formula
\[
 Z(z)=\frac{z}{z-1}-z\int_1^\infty\{t\}t^{-z-1}\,dt
\]
gives $|Z'(z)|\le1+2+16\cdot4<128$ along the horizontal segment from
$\rho$ to $s$. Here $|z-1|\ge1$, the integral is bounded by
$1/\Re z\le2$, and the differentiated integral by $1/(\Re z)^2\le4$.
Since $Z(\rho)=0$, we have $|Z(s)|\le128 e$. In particular
\begin{equation}\label{eq:small-omega-product}
        |Z(s)D_B(s)|\le 2048(2Ae+a)\le 1/8.
\end{equation}

We next show from the same finite hypothesis that this product is close
to one. Let $u_j$ be the Taylor coefficients at $z_0$ of
$Z(z)D_B(z)-1$. Two estimates at different radii are needed. At radius
$1/2$ the ordinary Dirichlet series converge absolutely, and the full
reciprocal identity gives
\[
 Z(z)D_B(z)-1=-\sum_{m\ge1}\sum_{d>B}\mu(d)(md)^{-z}.
\]
The absolute Taylor coefficients of $t^{-z}$ about $z_0$, weighted by
$(1/2)^j$, sum to $t^{-\Re z_0+1/2}$. As $\Re z_0-1/2\ge2$, comparison
with $\sum_m 1/m^2\le2$ and $\sum_{d>B}1/d^2\le1/B$ gives
\begin{equation}\label{eq:omega-small-radius}
                       \sum_{j\ge0}|u_j|2^{-j}\le 2/B.
\end{equation}
The estimates follow from $1/d^2\le1/(d-1)-1/d$ for $d>1$.

At radius $R=2+e/2$, use the two integral representations instead.
For any of the functions $f$ here with $|f(t)|\le C t^\beta$ the absolute
Taylor coefficient of degree $j$ in $\int_1^\infty f(t)t^{-z-1}dt$
is bounded by $C/(\Re z_0-\beta)^{j+1}$. Expand the exponential on
bounded intervals and integrate $u^je^{-du}/j!$ after $u=\log t$,
$d=\Re z_0-\beta$, as in Section~\ref{sec:growth}. Summing the
coefficient majorants with weight $R^j$ gives $C/(\Re z_0-\beta-R)$.
For $D_B$ take $C=A$, $\beta=1/2$, and multiply by
$|z_0|+R\le20$, obtaining $40A/e$. For $Z$, the fractional-part
integral has $C=1$, $\beta=0$, giving a weighted bound of at most $40$
after multiplication by $z$. The term $z/(z-1)$ contributes at most
$(|z_0|+R)/(|z_0-1|-R)\le2$ by its geometric series. Thus the bound
for $Z$ is at most $64$. The sum of weighted absolute coefficients
of a product is at most the product of the two bounds, so
\begin{equation}\label{eq:omega-large-radius}
                       \sum_{j\ge0}|u_j|R^j\le 1+64\cdot40A/e\le V.
\end{equation}
The coefficients in the two estimates agree term by term. In a neighborhood
of $z_0$, the integral for $D_B$ is the finite sum already evaluated,
and the Euler formula for $Z$ is its ordinary Dirichlet series.
Differentiating there at $z_0$ computes the coefficients in either
representation. All derivative orders through any specified $j$ use the
moment bounds just given, with geometric remainders on a smaller radius.

Evaluate the expansion at $s=z_0-2$. By
\eqref{eq:omega-small-radius} the terms through degree $J$ give absolute
sum at most $4^J(2/B)\le1/8$. The remainder is bounded by
\eqref{eq:omega-large-radius} by $V(2/R)^{J+1}$. Bernoulli's inequality
gives $(R/2)^{J+1}=(1+e/4)^{J+1}\ge1+(J+1)e/4\ge8V$.
It follows that $|Z(s)D_B(s)-1|\le1/4$, contradicting
\eqref{eq:small-omega-product} with an absolute gap of at least $5/8$.

The predicate tested at each $n$ in the finite interval is just the strict
integer inequality in the theorem. A bounded scan either returns its first
true instance or certifies the opposite weak inequality throughout the
interval. The contradiction excludes the second output. All truncation,
precision and search bounds are primitive recursive in $A$ and the
requested precision. For example the geometric remainder at the evaluation
point is controlled by the same Bernoulli bound with $J$ enlarged by a
factor $2^j$ for $j$ extra bits. Thus this reasoning gives the positive
witness, with the stated primitive recursive bound.
\end{proof}
